\documentclass[11pt,a4paper]{article}
\usepackage[T1]{fontenc}
\usepackage[utf8]{inputenc}
\usepackage{newtxtext}
\usepackage{amsmath,amsthm,mathtools}
\usepackage{newtxmath}
\usepackage[margin=27mm,headheight=15pt]{geometry}
\usepackage{microtype,booktabs,tabularx,array}
\usepackage[dvipsnames]{xcolor}
\usepackage{enumitem,etoolbox,fancyhdr}
\usepackage{hyperref}
\hypersetup{colorlinks=true,linkcolor=MidnightBlue,citecolor=MidnightBlue,
urlcolor=MidnightBlue,pdftitle={Global Fueter Primitives and Spherical Residues},
pdfauthor={OpenAI},pdfsubject={Two-period obstruction theory, critical spherical singularities, and exact energy asymptotics}}
\usepackage{bookmark}
\urlstyle{same}
\numberwithin{equation}{section}
\newtheorem{theorem}{Theorem}[section]
\newtheorem{proposition}[theorem]{Proposition}
\newtheorem{lemma}[theorem]{Lemma}
\newtheorem{corollary}[theorem]{Corollary}
\theoremstyle{definition}
\newtheorem{definition}[theorem]{Definition}
\newtheorem{example}[theorem]{Example}
\theoremstyle{remark}
\newtheorem{remark}[theorem]{Remark}
\newcommand{\R}{\mathbb R}
\newcommand{\C}{\mathbb C}
\newcommand{\HH}{\mathbb H}
\newcommand{\HC}{\mathbb H_{\C}}
\newcommand{\SSS}{\mathbb S}
\newcommand{\SR}{\mathcal{SR}}
\newcommand{\AM}{\mathcal{AM}}
\newcommand{\Aff}{\operatorname{Aff}_{\HH}}
\newcommand{\Fu}{\Delta_4}
\newcommand{\Dir}{\mathcal D}
\newcommand{\Per}{\operatorname{Per}}
\newcommand{\Ind}{\operatorname{Ind}}
\newcommand{\Res}{\operatorname{Res}}
\newcommand{\ReC}{\operatorname{Re}_{\C}}
\newcommand{\ImC}{\operatorname{Im}_{\C}}
\newcommand{\ImH}{\operatorname{Im}_{\HH}}
\newcommand{\ReH}{\operatorname{Re}_{\HH}}
\newcommand{\dd}{\,\mathrm d}
\newcommand{\dist}{\operatorname{dist}}
\newcommand{\ii}{\mathbf i}
\newcommand{\jj}{\mathbf j}
\newcommand{\kk}{\mathbf k}
\newcommand{\ax}{\mathrm{ax}}
\newcommand{\norm}[1]{\lVert#1\rVert}
\newcommand{\abs}[1]{\lvert#1\rvert}
\setlist{itemsep=3pt,topsep=5pt}
\setlength{\parskip}{2pt}
\setlength{\parindent}{1.15em}
\setlength{\emergencystretch}{2em}
\pagestyle{fancy}
\fancyhf{}
\fancyhead[L]{\small Global Fueter primitives and spherical residues}
\fancyhead[R]{\small\thepage}
\renewcommand{\headrulewidth}{0.3pt}
\title{\vspace{-10mm}\textbf{Global Fueter Primitives}\\[2pt]\textbf{and Spherical Residues}\\[6pt]
\large A two-period obstruction theory, critical-growth classification,\\
and exact energy asymptotics}
\author{Research manuscript}
\date{21 September 2026}

\begin{document}
\maketitle
\thispagestyle{empty}
\begin{abstract}
The local inverse Fueter theorem does not by itself decide whether an axial
monogenic function has a single-valued slice-regular primitive on a domain
with holes. We give a complete global criterion in terms of two closed
quaternion-valued one-forms computed directly from the target function.
Their periods are exactly the two coefficients of the affine monodromy of
local primitives. Explicit logarithmic modes realize every period assignment
on a finitely connected planar domain, yielding a split cokernel of real
dimension eight per hole, together with quantitative non-approximation and
continuous inversion estimates. A holomorphic second-derivative invariant
converts the obstruction into two ordinary complex moments, or a constrained Laurent
coefficient pair at an isolated puncture.

The same periods classify every axial monogenic singularity along a
conjugacy sphere whose growth is at most the reciprocal of the distance to
that sphere. The singular part has precisely two quaternionic coefficients;
the remainder extends regularly. We prove the sharp removability threshold,
an exact leading supremum formula, and a logarithmic energy law with an
explicit positive-definite residue quadratic form. We also obtain a finite
spherical principal-part theorem with polynomial growth and far-field moment
conditions. The foundational inversion theory and the special logarithmic
kernels are credited to earlier work. The global and quantitative formulations
proved here are the proposed contributions; their priority is qualified by the
literature assessment in the text.
\end{abstract}
\medskip
\noindent\textbf{Keywords.} Quaternionic analysis; slice regularity; inverse
Fueter theorem; periods; monodromy; spherical singularities; residues.

\medskip
\noindent\textbf{Status.} All asserted mathematical results below have proofs.
This is a research manuscript, not a peer-reviewed priority claim. No assertion
is made that a named, previously published open conjecture has been resolved.
The claimed contribution is a coherent collection of theorems not located in
the primary sources inspected, with the strongest new formulations identified
in \S\ref{sec:position}. Symbolic and numerical checks accompany, but do not
replace, the proofs.

\newpage
\begingroup\small
\setlength{\parskip}{0pt}
\makeatletter
\patchcmd{\l@section}{1.0em}{0.45em}{}{}
\makeatother
\tableofcontents
\endgroup
\newpage

\section{The global problem and the main results}
\label{sec:intro}
The Fueter map in quaternionic dimension four is the ordinary Laplacian
\[
 \Fu:\SR\longrightarrow\AM,
\]
from left slice-regular functions to left axial Cauchy--Fueter-regular
functions. Its kernel consists of the right-affine functions $q\mapsto qa+b$
with $a,b\in\HH$. Local inversion therefore has an eight-real-dimensional
ambiguity. The question addressed here is what happens when local inverses
are continued around a hole.

The supplied expository manuscript \cite{companion} ends its constructive
inverse theorem with precisely this issue: two successive closed one-forms
can acquire periods on a non-simply-connected domain. We replace that
sequential construction by two \emph{simultaneously defined} closed forms.
This makes the global obstruction measurable without first choosing either
primitive. More importantly, it connects the obstruction to the complete
critical-order singular part and to an exact energy coefficient.

\subsection{A compact statement of the results}
Let $U$ be a connected open subset of the upper complex half-plane, with
coordinates $z=x+\iota r$, $r>0$. Its circularization is
\[
 \Omega_U=\{x+Ir:(x,r)\in U,\ I\in\SSS\},
 \qquad \SSS=\{I\in\HH:I^2=-1\}.
\]
Suppose $g(x+Ir)=P(x,r)+IQ(x,r)$ is axial monogenic. Define
\begin{align}
 \omega_g&=\frac12(-P\dd x+Q\dd r),\label{eq:intro-omega}\\
 \beta_g&=\frac12\bigl((xP+rQ)\dd x+(rP-xQ)\dd r\bigr).
 \label{eq:intro-beta}
\end{align}
The central assertions, proved in the following sections, are these.

\begin{enumerate}[label=\textup{(\roman*)}]
\item A global single-valued slice-regular $f$ with $\Fu f=g$ exists
exactly when both forms have zero periods on all closed curves in $U$.
Around a loop $\gamma$, a local primitive acquires precisely
$q a_\gamma+b_\gamma$, where
$a_\gamma=\int_\gamma\omega_g$ and $b_\gamma=\int_\gamma\beta_g$.
\item If $U$ has a finite winding basis of $m$ holes, then
$\AM(\Omega_U)/\Fu\SR(\Omega_U)$ has quaternionic dimension $2m$.
We construct an explicit continuous linear splitting. Real punctures do not
create additional obstructions when the domain meets the real axis.
\item A globally defined holomorphic $\HC$-valued function
\[
 H_g=-\frac12(P+rP_r)-\frac{\iota r}{2}P_x
\]
satisfies $a_\gamma=\int_\gamma H_g\dd z$ and
$b_\gamma=-\int_\gamma zH_g\dd z$. Thus two complex moments, not just one
residue, are relevant.
\item Near a sphere $S_p=s+t\SSS$, every axial monogenic $g$ satisfying
$|g(q)|=O(\dist(q,S_p)^{-1})$ is uniquely
\[
 g=G_{p;a,b}+h,
\]
where $h$ extends regularly across $S_p$ and $(a,b)$ are its two periods.
Writing
\begin{equation}
 \mathcal R_p(a,b)^2=|sa+b|^2+t^2|a|^2,\label{eq:intro-resnorm}
\end{equation}
the energy in a thin tube has the asymptotic
\[
 \int_{T_{\rho_0}(S_p)\setminus\overline{T_\rho(S_p)}}|g|^2\dd V
 =8\mathcal R_p(a,b)^2\log\frac{\rho_0}{\rho}
   +E_{\mathrm{ren}}+o(1).
\]
In particular, $o(\dist^{-1})$ growth is removable, whereas every nonzero
residue pair produces a nonremovable critical singularity.
\end{enumerate}

These are not merely rephrasings of local solvability. They give the exact
size of the global failure of surjectivity, a computable obstruction basis,
and a classification with a sharp quantitative invariant. Ordinary
single-point Fueter residues in four-dimensional space are different objects:
here the singular set is an entire two-sphere.

\subsection{Conventions and prerequisites}
We use classical one-variable holomorphic function theory, elementary
integration of closed one-forms, the winding number, and the standard
mean-value and interior derivative estimates for harmonic functions.
Quaternionic analytic statements needed in the proofs are derived below.
The basic slice-function framework is that of Ghiloni and Perotti
\cite{GP}. All vector spaces over $\HH$ in this paper are \emph{right}
vector spaces. An auxiliary complex scalar is central; it must never be
confused with a quaternionic imaginary unit.

\section{Analytic framework and local reconstruction}
\label{sec:framework}
\subsection{Stems and axial functions}
Write $q=x_0+\ii x_1+\jj x_2+\kk x_3$, with
$\ii^2=\jj^2=\kk^2=\ii\jj\kk=-1$.
Quaternionic conjugation is denoted by a bar and reverses products.
The norm is Euclidean and multiplicative. Put
\[
 \HC=\{A+\iota B:A,B\in\HH\},\qquad \iota^2=-1,
\]
where $\iota$ commutes with $\HH$. The operations
$\ReC(A+\iota B)=A$ and $\ImC(A+\iota B)=B$ take values in $\HH$.
We use the norm
\[
 |A+\iota B|_{\HC}^2=|A|^2+|B|^2.
\]
Multiplication by a central complex scalar of modulus $\lambda$ multiplies
this norm by $\lambda$.

Let $F=A+\iota B:U\to\HC$ be holomorphic, equivalently
\begin{equation}
 A_x=B_r,\qquad A_r=-B_x.\label{eq:CR}
\end{equation}
It induces
\begin{equation}
 f(x+Ir)=A(x,r)+IB(x,r).\label{eq:induced}
\end{equation}
We call such $f$ slice regular and write $f\in\SR(\Omega_U)$.
Equivalently, reflect $F$ to $U\cup\overline U$ by
$F(\bar z)=\kappa(F(z))$, where $\kappa(A+\iota B)=A-\iota B$.
This is a product-domain definition; the domain need not meet $\R$.
It rules out unrelated slice-wise functions not induced by one stem.

The Cauchy--Fueter operator is
\[
 \Dir=\partial_{x_0}+\ii\partial_{x_1}
                   +\jj\partial_{x_2}+\kk\partial_{x_3},
 \qquad \Fu=\sum_{\mu=0}^3\partial_{x_\mu}^2.
\]
An axial function is $g(x+Ir)=P(x,r)+IQ(x,r)$, with quaternion-valued
$P,Q$ independent of $I$. We write $\AM(\Omega_U)$ for the smooth axial
functions satisfying $\Dir g=0$. Requiring all $I$ makes $P,Q$ unique:
the values at $I$ and $-I$ recover them.

\begin{lemma}[Axial differential identities]\label{lem:axial}
For smooth axial functions on $\Omega_U$,
\begin{align}
 \Dir(A+IB)&=A_x-B_r-2B/r+I(B_x+A_r),\label{eq:Dir-axial}\\
 \Fu(A+IB)&=A_{xx}+A_{rr}+2A_r/r
 +I(B_{xx}+B_{rr}+2B_r/r-2B/r^2).\label{eq:lap-axial}
\end{align}
Consequently, $g=P+IQ$ belongs to $\AM(\Omega_U)$ exactly when
\begin{equation}
 P_x-Q_r=2Q/r,\qquad Q_x+P_r=0.\label{eq:vekua}
\end{equation}
If $f=A+IB$ is slice regular, then
\begin{equation}
 \Fu f=\frac{2A_r}{r}
       +I\left(\frac{2B_r}{r}-\frac{2B}{r^2}\right).
 \label{eq:fueter}
\end{equation}
This function is axial monogenic.
\end{lemma}
\begin{proof}
Put $v=\ii x_1+\jj x_2+\kk x_3$, $r=|v|$, and $I=v/r$.
For $\ell=1,2,3$, writing $e_1=\ii,e_2=\jj,e_3=\kk$, we have
\[
 \partial_{x_\ell}r=x_\ell/r,\qquad
 \partial_{x_\ell}I=e_\ell/r-vx_\ell/r^3.
\]
Thus $\sum e_\ell\partial_{x_\ell}I=-2/r$,
$\Delta_{\R^3}I=-2I/r^2$, and
$\sum(\partial_{x_\ell}I)(\partial_{x_\ell}r)=0$.
The product rule gives \eqref{eq:Dir-axial}--\eqref{eq:lap-axial}.
Comparing the values at $I$ and $-I$ gives \eqref{eq:vekua}.

The Cauchy--Riemann equations make $A$ and $B$ planar harmonic, reducing
the Laplacian to \eqref{eq:fueter}. Set
$P=2A_r/r$ and $Q=2B_r/r-2B/r^2$. Direct differentiation using
\eqref{eq:CR} yields
\[
 P_x=2B_{rr}/r=Q_r+2Q/r,\qquad Q_x+P_r=0.
\]
Hence \eqref{eq:vekua} holds.
\end{proof}
The mapping assertion is the classical quaternionic Fueter theorem, not a
new claim. The formulas fix the normalization used in all periods and energy
constants below. In particular, $\Fu(q^2)=-4$, and there is no hidden
factor of $1/2$ in the definition of the map.

\subsection{Two simultaneous potentials}
\begin{lemma}[Closed forms and potential identities]\label{lem:forms}
For $g=P+IQ\in\AM(\Omega_U)$, the forms
\eqref{eq:intro-omega} and \eqref{eq:intro-beta} are closed.
If $f=A+IB\in\SR(\Omega_U)$ satisfies $\Fu f=g$, then
\begin{equation}
 C=B/r,\qquad E=A-xB/r\label{eq:CE}
\end{equation}
satisfy
\begin{equation}
 dC=\omega_g,\qquad dE=\beta_g,\qquad f(q)=qC+E.
 \label{eq:potential-identities}
\end{equation}
Conversely, any primitives $C,E$ of the two forms induce a slice-regular
solution through $A=xC+E$, $B=rC$.
\end{lemma}
\begin{proof}
The coefficient of $dx\wedge dr$ in $d\omega_g$ is
$(Q_x+P_r)/2$, so it vanishes. For the second form, the equality of mixed
coefficients is the identity
\begin{align*}
 &\partial_r(xP+rQ)-\partial_x(rP-xQ)\\
 &\hspace{15mm}=x(P_r+Q_x)+2Q+r(Q_r-P_x)=0.
\end{align*}
All computations are componentwise in the real vector space $\HH$.

For a primitive $f$, \eqref{eq:CR} and \eqref{eq:fueter} give
\[
 C_x=B_x/r=-P/2,\qquad C_r=B_r/r-B/r^2=Q/2.
\]
Moreover,
\begin{align*}
 E_x&=A_x-C-xC_x
     =B_r-C+xP/2=(rQ+xP)/2,\\
 E_r&=A_r-xC_r=(rP-xQ)/2.
\end{align*}
This proves \eqref{eq:potential-identities}.

Conversely, the primitive equations imply
\begin{align*}
 A_x&=C+x(-P/2)+(xP+rQ)/2=C+rQ/2=B_r,\\
 A_r&=xQ/2+(rP-xQ)/2=rP/2=-B_x.
\end{align*}
Thus $A+\iota B$ is holomorphic, and substituting in
\eqref{eq:fueter} yields $\Fu f=P+IQ$.
\end{proof}

The important change of variable is $E=A-xC$. The sequential local inverse
would first integrate $dC=\omega_g$ and then integrate
\[
 dA=(C+rQ/2)\dd x+(rP/2)\dd r.
\]
Subtracting $d(xC)$ removes the unknown $C$ and gives $\beta_g$ directly.
No regular product, reciprocal, or noncommutative logarithm is being used.

\begin{corollary}[Local inverse and precise kernel]\label{cor:local}
Every $g\in\AM(\Omega_U)$ has a slice-regular Fueter primitive on the
circularization of each simply connected subdomain of $U$.
On connected $U$,
\begin{equation}
 \ker(\Fu|_{\SR})=\Aff
 =\{q\mapsto qa+b:a,b\in\HH\}.\label{eq:kernel}
\end{equation}
\end{corollary}
\begin{proof}
Integrate the two closed forms on a simply connected set and apply
Lemma~\ref{lem:forms}. If $\Fu f=0$, then $dC=dE=0$, so connectedness
makes both potentials constant. The converse follows immediately from
\eqref{eq:fueter}.
\end{proof}
This reproves the needed local portion of inverse Fueter theory
\cite{CSS,CPSS,Perotti}. Notice that the constants in \eqref{eq:kernel}
are quaternions, not arbitrary elements of $\HC$. A complex-affine stem
with nonreal central coefficients need not be in the Fueter kernel.

\section{The two-period theorem and affine monodromy}
\label{sec:period}
For a closed piecewise $C^1$ curve $\gamma$ in $U$, put
\begin{equation}
 a_\gamma(g)=\int_\gamma\omega_g,\qquad
 b_\gamma(g)=\int_\gamma\beta_g.\label{eq:periods}
\end{equation}
The integrals use the indicated order of quaternionic factors. Here the
coefficients of $dx,dr$ multiply real differentials, so no ordering ambiguity
remains. Since the forms are closed, the integrals depend only on the
integral homology class of $\gamma$.

\begin{theorem}[Global single-valued inversion]\label{thm:period}
For connected open $U\subset\{r>0\}$ and $g\in\AM(\Omega_U)$, the
following are equivalent:
\begin{enumerate}[label=\textup{(\roman*)}]
\item There is $f\in\SR(\Omega_U)$ with $\Fu f=g$.
\item Both $\omega_g$ and $\beta_g$ are exact on $U$.
\item $a_\gamma(g)=b_\gamma(g)=0$ for every closed piecewise $C^1$ curve
$\gamma\subset U$.
\end{enumerate}
When these conditions hold, choose $z_0\in U$ and constants $c_0,e_0\in\HH$.
Every solution is obtained from
\begin{equation}
 C(z)=c_0+\int_{z_0}^{z}\omega_g,\qquad
 E(z)=e_0+\int_{z_0}^{z}\beta_g,\qquad
 f(x+Ir)=(x+Ir)C(z)+E(z).\label{eq:global-inverse}
\end{equation}
The integrals are independent of the paths. The solution is unique modulo
$\Aff$.
\end{theorem}
\begin{proof}
Lemma~\ref{lem:forms} proves (i)$\Leftrightarrow$(ii). Exact forms have
zero periods, so (ii) implies (iii). For the reverse implication, an open
connected subset of the plane is path connected by polygonal paths.
Define the two integrals in \eqref{eq:global-inverse} along such paths.
Two choices differ by a closed path; condition (iii) makes both differences
zero. On a small disk the usual differentiation of a path integral shows
that these functions are primitives of the corresponding forms.
They therefore satisfy (ii) globally. Formula \eqref{eq:global-inverse}
and the uniqueness statement now follow from Lemma~\ref{lem:forms} and
Corollary~\ref{cor:local}.
\end{proof}

\begin{theorem}[Complete affine monodromy]\label{thm:monodromy}
A local Fueter primitive can be analytically continued along every path in
$U$. Continuation around a based loop $\gamma$ changes its stem and its
induced slice function by
\begin{equation}
 F\longmapsto F+z a_\gamma+b_\gamma,\qquad
 f\longmapsto f+q a_\gamma+b_\gamma.\label{eq:monodromy}
\end{equation}
Consequently the full monodromy is the additive homomorphism
\begin{equation}
 H_1(U;\mathbb Z)\longrightarrow\HH\oplus\HH,
 \qquad [\gamma]\longmapsto(a_\gamma,b_\gamma).\label{eq:monodromy-map}
\end{equation}
There are no further local compatibility obstructions.
\end{theorem}
\begin{proof}
Pull the closed forms to the universal covering surface of $U$, or,
equivalently, integrate them along paths from a fixed base point while
retaining the homotopy class of the path. They have global potentials on
the cover. Local complex coordinates are inherited from $U$, and
Lemma~\ref{lem:forms} makes the corresponding $F=zC+E$ holomorphic.
It agrees with the prescribed primitive after choosing the initial
constants correctly. Appending $\gamma$ adds $a_\gamma$ to $C$ and
$b_\gamma$ to $E$, which proves \eqref{eq:monodromy}.

Concatenation adds periods and reversal negates them. Integration of
closed forms is invariant under homology, so the homomorphism factors
through $H_1(U;\mathbb Z)$. Conversely, the construction supplies an
actual continued primitive for every path; no additional transition
condition remains.
\end{proof}

\begin{remark}[What is, and is not, noncommutative here]
The target algebra is noncommutative, but the obstruction group is additive
and abelian. It has two quaternionic coordinates because the kernel is
affine. A scalar complex residue alone cannot retain both coordinates.
The geometry of quaternionic conjugacy spheres will re-enter essentially
in the singularity norm and energy theorems.
\end{remark}

\section{Logarithmic modes and the split global obstruction}
\label{sec:modes}
\subsection{Branch-independent modes with prescribed periods}
Fix $p=s+\iota t$ with $t>0$. On a simply connected patch avoiding $p$ and
$\bar p$, choose a branch of
\begin{equation}
 L_p(z)=\frac1{2\pi\iota}\log\frac{z-p}{z-\bar p}.\label{eq:log-mode}
\end{equation}
For $a,b\in\HH$, define the local holomorphic stem
\begin{equation}
 F_{p;a,b}(z)=L_p(z)(za+b).\label{eq:mode-stem}
\end{equation}

\begin{proposition}[Global mode and period normalization]\label{prop:modes}
The local functions $\Fu\mathcal I(F_{p;a,b})$ agree on overlaps and define
a single-valued axial monogenic function $G_{p;a,b}$ on
$\HH\setminus S_p$, where $S_p=s+t\SSS$.
The function is smooth across the real axis. On every closed curve
$\gamma\subset\{r>0\}\setminus\{p\}$,
\begin{equation}
 a_\gamma(G_{p;a,b})=\Ind(\gamma,p)a,\qquad
 b_\gamma(G_{p;a,b})=\Ind(\gamma,p)b.\label{eq:mode-periods}
\end{equation}
\end{proposition}
\begin{proof}
Two branches of the logarithm differ locally by $2\pi\iota n$,
$n\in\mathbb Z$. Their $L_p$ differ by the real integer $n$; the resulting
stems differ by $n(za+b)$. This induces an element of $\Aff$, killed by
$\Fu$. Lemma~\ref{lem:axial} makes the common result axial monogenic.

At a real point $x$, the ratio $(x-p)/(x-\bar p)$ has modulus one and
is nonzero. In a small symmetric disk about $x$, choose an analytic
logarithm. Along its real diameter that logarithm is purely imaginary,
so $L_p$ is real there. Schwarz reflection gives the stem symmetry for
$L_p$, hence for \eqref{eq:mode-stem}. The induced slice-regular local
primitive is real analytic across the real axis. Its Laplacian supplies
a smooth continuation of $G_{p;a,b}$ there. Branch independence glues
these extensions to the already defined function.

Continuing the logarithm around $\gamma$ changes $L_p$ by
$\Ind(\gamma,p)-\Ind(\gamma,\bar p)$. The second index is zero because
the loop lies in the upper half-plane. Its affine monodromy is therefore
exactly $z\Ind(\gamma,p)a+\Ind(\gamma,p)b$.
Theorem~\ref{thm:monodromy} gives \eqref{eq:mode-periods}.
\end{proof}

\subsection{An explicit formula with no logarithm branches}
The imaginary central component of $L_p$ is single valued:
\begin{equation}
 v_p(x,r)=\ImC L_p(x+\iota r)
 =-\frac1{4\pi}\log
   \frac{(x-s)^2+(r-t)^2}{(x-s)^2+(r+t)^2}.\label{eq:v}
\end{equation}
For $r>0$, write $v=v_p$. The two components of the mode are
\begin{align}
 P_{p;a,b}&=-2(v_r+v/r)a-\frac{2v_x}{r}(xa+b),\label{eq:mode-P}\\
 Q_{p;a,b}&=-2v_xa+2\left(\frac{v_r}{r}-\frac{v}{r^2}\right)(xa+b).
 \label{eq:mode-Q}
\end{align}
Indeed, write a branch as $L_p=u+\iota v$. Then
\[
 A=u(xa+b)-rv a,\qquad B=rua+v(xa+b),
 \qquad C=u a+\frac vr(xa+b).
\]
Since $u_x=v_r$ and $u_r=-v_x$, the equations $P=-2C_x$ and $Q=2C_r$
give \eqref{eq:mode-P}--\eqref{eq:mode-Q}.
These are useful both for implementation and for seeing directly that the
multivalued real component $u$ has disappeared.

The apparent denominators at $r=0$ are removable. More explicitly, putting
$d=(x-s)^2+t^2$, Taylor expansion of \eqref{eq:v} gives
$v=t r/(\pi d)+O(r^3)$ and hence
\begin{equation}
 G_{p;a,b}(x)=-\frac{4t}{\pi d}a
       +\frac{4t(x-s)}{\pi d^2}(xa+b),\qquad x\in\R.
 \label{eq:mode-real}
\end{equation}
The $Q$ component is $O(r)$. Smoothness of every order also follows from
the local stem proof in Proposition~\ref{prop:modes}.

\begin{remark}[The modes themselves have predecessors]\label{rem:known-modes}
For $p=\iota$,
\[
 L'_\iota(z)=\frac1{\pi(1+z^2)},
\]
so locally $L_\iota=\pi^{-1}\arctan z+c$ with real $c$ on a reflected
patch. The functions with local primitives
$(2\pi)^{-1}\arctan z$ and $(2\pi)^{-1}z\arctan z$ are the spherical
Cauchy kernels $N^+$ and $N^-$ studied in
\cite[Example~2]{CPSS}. Thus our $G_{\iota;0,1}$ and
$G_{\iota;1,0}$ are respectively $2N^+$ and $2N^-$ in that
normalization. We do not claim the invention of these kernels. Their
period normalization, use as a complete obstruction basis, and the
classification and energy laws below are the formulations developed here.
\end{remark}

\subsection{Finite topology and exact splitting}
We make the modest topological hypothesis explicit rather than imposing
unnecessary boundary regularity on $U$.

\begin{definition}[Finite winding basis]\label{def:winding}
A connected plane domain $U\subset\{r>0\}$ has a finite winding basis
of size $m$ if there are piecewise $C^1$ closed curves
$\gamma_1,\ldots,\gamma_m$ in $U$ and points
$p_1,\ldots,p_m\in\{r>0\}\setminus U$ such that their homology classes
form a basis of $H_1(U;\mathbb Z)$ and
\begin{equation}
 \Ind(\gamma_j,p_k)=\delta_{jk}.\label{eq:winding-basis}
\end{equation}
The case $m=0$ means $H_1(U;\mathbb Z)=0$.
\end{definition}
For example, the upper half-plane with $m$ distinct points removed has
this property: take small positively oriented circles and the removed
points. The usual finitely connected domains obtained by removing $m$
disjoint closed Jordan regions from a simply connected domain also admit
such a basis, with one point in each removed region. The abstract theorem
below uses only Definition~\ref{def:winding}.

\begin{theorem}[Split cokernel and complete obstruction dimension]
\label{thm:split}
If $U$ has a finite winding basis of size $m$, define
\[
 \Per(g)=\bigl((a_{\gamma_j}(g),b_{\gamma_j}(g))\bigr)_{j=1}^m.
\]
There is an exact sequence of right quaternionic vector spaces
\begin{equation}
 0\longrightarrow\Aff\longrightarrow\SR(\Omega_U)
 \xrightarrow{\ \Fu\ }\AM(\Omega_U)
 \xrightarrow{\ \Per\ }\HH^{2m}\longrightarrow0.\label{eq:exact-sequence}
\end{equation}
The last surjection has the explicit section
\begin{equation}
 \mathcal S\bigl((a_j,b_j)_{j=1}^m\bigr)
       =\sum_{j=1}^mG_{p_j;a_j,b_j}.\label{eq:section}
\end{equation}
Consequently
\begin{equation}
 \AM(\Omega_U)=\Fu\SR(\Omega_U)\ \oplus\ \mathcal S(\HH^{2m}),
 \qquad
 \dim_{\R}\frac{\AM(\Omega_U)}{\Fu\SR(\Omega_U)}=8m.
 \label{eq:direct-sum}
\end{equation}
Every $g$ has the unique decomposition
\begin{equation}
 g=\Fu f+\sum_{j=1}^mG_{p_j;a_{\gamma_j}(g),b_{\gamma_j}(g)},
 \label{eq:decomposition}
\end{equation}
where $f$ is unique modulo $\Aff$.
\end{theorem}
\begin{proof}
The first kernel is \eqref{eq:kernel}. Vanishing on the homology basis is
vanishing of all periods, so Theorem~\ref{thm:period} gives
$\ker\Per=\Fu\SR$. Proposition~\ref{prop:modes} and
\eqref{eq:winding-basis} give $\Per\mathcal S=\operatorname{id}$.
This proves surjectivity, exactness, and splitting. Subtracting
$\mathcal S\Per(g)$ leaves zero periods and hence a global image of
$\Fu$. Applying $\Per$ to any two such decompositions proves uniqueness
of the mode part; the affine kernel gives the remaining ambiguity.
There are $2m$ free quaternionic coordinates, hence $8m$ real coordinates.
\end{proof}

\begin{example}[The second obstruction cannot be omitted]\label{ex:second}
Let $U=\{r>0\}\setminus\{\iota\}$ and let $\gamma$ be a positive small
circle about $\iota$. The function $G_{\iota;0,1}$ has
\[
 \int_\gamma\omega=0,\qquad\int_\gamma\beta=1.
\]
The first potential $C$ therefore exists globally, but no global
slice-regular Fueter primitive exists. Conversely $G_{\iota;1,0}$ has
periods $(1,0)$. The two failure mechanisms are independent, even for
modes whose stem components are real valued.
\end{example}

\section{Quantitative stability and non-approximation}
\label{sec:stability}
The obstruction is not destroyed by a small perturbation on its detecting
contours. For a compact $K\subset U$, define
\begin{equation}
 \norm{g}_{\ax,K}=\sup_{z\in K}\sqrt{|P(z)|^2+|Q(z)|^2}.
 \label{eq:axnorm}
\end{equation}
This is equivalent to the usual supremum norm on the circularization:
\begin{equation}
 \norm{g}_{\ax,K}\le
 \sup_{z\in K,I\in\SSS}|P(z)+IQ(z)|
 \le\sqrt2\norm{g}_{\ax,K}.\label{eq:norm-equivalence}
\end{equation}
For the first inequality, average the squared norm over $I$; the cross
term integrates to zero. The second follows from the triangle inequality.

\begin{proposition}[Period bounds and a lower bound for approximation]
\label{prop:distance}
Let $\gamma$ have length $\ell_\gamma$ and put
$R_\gamma=\max_{z\in\gamma}|z|>0$. Then
\begin{equation}
 |a_\gamma(g)|\le\frac{\ell_\gamma}{2}\norm{g}_{\ax,\gamma},
 \qquad
 |b_\gamma(g)|\le\frac{R_\gamma\ell_\gamma}{2}
                                  \norm{g}_{\ax,\gamma}.
 \label{eq:period-bounds}
\end{equation}
For every global $f\in\SR(\Omega_U)$,
\begin{equation}
 \norm{g-\Fu f}_{\ax,\gamma}\ge
 \frac{2}{\ell_\gamma}
 \max\left\{|a_\gamma(g)|,\frac{|b_\gamma(g)|}{R_\gamma}\right\}.
 \label{eq:distance-lower}
\end{equation}
\end{proposition}
\begin{proof}
For a tangent vector $(u,v)$, Cauchy--Schwarz in $\HH\oplus\HH$ gives
\[
 |\omega_g(u,v)|\le\tfrac12\sqrt{|P|^2+|Q|^2}\sqrt{u^2+v^2}.
\]
The real matrix
$\left(\begin{smallmatrix}x&r\\r&-x\end{smallmatrix}\right)$
multiplies Euclidean lengths by $|z|$. Applying the same estimate to
$\beta_g$ gives the second bound in \eqref{eq:period-bounds} after
integration. By Theorem~\ref{thm:period}, the periods of $\Fu f$ vanish.
Apply both bounds to $g-\Fu f$ and rearrange.
\end{proof}
Thus a target with a nonzero period cannot even be approximated uniformly
on a detecting loop by Fueter images of globally single-valued primitives.
This is stronger than merely saying that no exact global solution exists.

\begin{theorem}[Continuous projection and normalized inverse]
\label{thm:continuous}
For a finite winding basis, the projection
\begin{equation}
 \Pi=\operatorname{id}-\mathcal S\Per
 \label{eq:projection}
\end{equation}
is continuous in the compact-open smooth topology on $\AM(\Omega_U)$,
has range $\Fu\SR(\Omega_U)$, and has kernel $\mathcal S(\HH^{2m})$.
The range is closed even for compact-open convergence without derivatives.

Fix $z_0\in U$. On the zero-period subspace, there is a unique inverse
$\mathcal T g=f$ normalized by $C(z_0)=E(z_0)=0$ in \eqref{eq:CE}.
For every compact $K\subset U$, there are a compact $K_*\subset U$
containing $K\cup\{z_0\}$ and a finite $L_K$ such that
\begin{equation}
 \sup_{z\in K,I\in\SSS}|\mathcal T g(x+Ir)|
 \le L_K R_*\norm{g}_{\ax,K_*},\qquad
 R_*=\max_{z\in K_*}|z|.\label{eq:inverse-estimate}
\end{equation}
The normalized inverse is continuous in the compact-open smooth topology.
\end{theorem}
\begin{proof}
Each period is a continuous linear functional by
\eqref{eq:period-bounds}. A finite sum of fixed smooth modes with these
coefficients depends continuously on $g$. The algebraic identities follow
from $\Per\mathcal S=\operatorname{id}$. The kernel of the finite
continuous period map is closed under compact-open convergence, proving
the assertion about the range within $\AM$.

To obtain a controlled family of paths, cover $K$ by finitely many small
open disks with compact closures in $U$. Join $z_0$ to each disk center
by a fixed polygonal path. A final straight segment reaches any point
of the disk. The union of these paths and closed disks is a compact
$K_*$, and their lengths have a common bound $L_K$.

Integrating \eqref{eq:period-bounds} in its pointwise form along these
paths gives
\[
 |C(z)|\le\tfrac12L_K\norm{g}_{\ax,K_*},\qquad
 |E(z)|\le\tfrac12R_*L_K\norm{g}_{\ax,K_*}.
\]
Since $|x+Ir|=|z|\le R_*$, the identity $f=qC+E$ proves
\eqref{eq:inverse-estimate}. Differentiating $dC=\omega_g$ and
$dE=\beta_g$ shows that every positive-order derivative of $C,E$ is
controlled locally by finitely many derivatives of $P,Q$. Together with
the zeroth-order estimate and the fact that compact subsets of $U$ stay
away from $r=0$, this proves smooth continuity of the induced inverse.
\end{proof}
The normalization sets the primitive to zero on the entire base sphere,
not only at a chosen point. This uses precisely the two quaternionic
constants available in the affine ambiguity. The constants in
\eqref{eq:inverse-estimate} are domain- and path-dependent; no uniform
estimate for degenerating domains is claimed.

\section{Domains meeting the real axis}
\label{sec:realaxis}
Let $D\subset\C$ be connected, open, invariant under complex conjugation,
and meeting $\R$. Its circularization is
\[
 \Omega_D=\{x+Iy:x+\iota y\in D,\ I\in\SSS\}.
\]
A holomorphic stem on $D$ satisfies
$F(\bar z)=\kappa(F(z))$. Equivalently, its components $A,B$ are even
and odd in $y$, respectively. This definition produces the usual
slice-regular functions on an axially symmetric slice domain.

For a smooth axial $g$ on $\Omega_D$, the signed components $P,Q$ are
smooth on $D$, with $P$ even and $Q$ odd in $y$. For example, fixing
$I\in\SSS$, they are recovered by
\[
 P(x,y)=\frac{g(x+Iy)+g(x-Iy)}2,\qquad
 Q(x,y)=-\frac I2\bigl(g(x+Iy)-g(x-Iy)\bigr).
\]
This also explains the parity and smoothness assertions at $y=0$.
We continue to write $\AM(\Omega_D)$ for this axial monogenic class.

\begin{lemma}[The upper half of a symmetric domain]\label{lem:upper}
If $D$ is as above, $U=D\cap\{\operatorname{Im}z>0\}$ is connected.
\end{lemma}
\begin{proof}
Join any two points of $U$ by a path in $D$. Reflect each portion below
the real axis upward by replacing $x+\iota y$ by $x+\iota|y|$.
The resulting continuous path lies in $D\cap\{y\ge0\}$ because $D$
is symmetric. Its compact image has a positive distance from the closed
complement of $D$. Translating the entire path upward by a sufficiently
small positive imaginary constant keeps it inside $D$ and makes it lie
strictly above $\R$. Short vertical segments at the two endpoints, also
contained in $D\cap\{y>0\}$, complete a path in $U$.
\end{proof}

\begin{theorem}[The upper-period criterion across the real axis]
\label{thm:realaxis}
For $g\in\AM(\Omega_D)$, a global $f\in\SR(\Omega_D)$ with $\Fu f=g$
exists if and only if both forms \eqref{eq:intro-omega} and
\eqref{eq:intro-beta}, restricted to $U=D\cap\{y>0\}$, have zero
periods there. The primitive is unique modulo $\Aff$.
If $U$ has a finite winding basis of size $m$, the exact sequence and
splitting of Theorem~\ref{thm:split} hold with $\Omega_D$ in place of
$\Omega_U$ and with these upper periods.
\end{theorem}
\begin{proof}
Necessity follows by restriction and Theorem~\ref{thm:period}.
For sufficiency, obtain potentials $C,E$ on connected $U$ from that
theorem. Replace $r$ by the signed coordinate $y$ in the two forms.
They are smooth on $D$ and closed even at $y=0$: the closedness identities
hold away from the axis and extend to it by continuity.

Reflection $\tau(x,y)=(x,-y)$ leaves both forms invariant. Indeed,
$P$ is even and $Q$ odd; the $dx$ coefficients in both forms are even,
and their $dy$ coefficients are odd. Take a small symmetric disk about
a real point of $D$. Each form has a local primitive on the disk.
For such a primitive $C_0$, reflection invariance gives
$d(C_0\circ\tau-C_0)=0$. This difference vanishes on the real diameter,
so it is zero; $C_0$ is even. The same applies to a local primitive $E_0$.
On the connected upper half of the disk, $C_0$ differs from $C$ by a
constant, and $E_0$ differs from $E$ by a constant. Adjust those constants
and extend $C,E$ across the diameter. All such local extensions agree on
overlaps, since they agree above the axis and are even. Finally reflect
the upper functions to the rest of the lower domain.

The resulting global $C,E$ are smooth and even. Therefore
$A=xC+E$ is even and $B=yC$ is odd. The Cauchy--Riemann calculation in
Lemma~\ref{lem:forms} holds off the axis and extends to it by continuity.
Thus $A+\iota B$ is a holomorphic stem on $D$. Its induced $f$ solves
$\Fu f=g$, including on the real axis by continuity.

Uniqueness follows by restriction to $U$ and reflection. For the finite
basis statement, each omitted upper point $p_j$ has its entire sphere
outside $\Omega_D$, and Proposition~\ref{prop:modes} supplies a mode
smooth across the real axis. Hence the same section realizes all upper
periods, and the preceding criterion proves exactness.
\end{proof}

\begin{corollary}[Deleted spheres and deleted real points]
\label{cor:deleted}
Let $p_1,\ldots,p_m$ be distinct upper-half-plane points and
$c_1,\ldots,c_\ell$ distinct real numbers. On
\begin{equation}
 \Omega=\HH\setminus\left(\bigcup_{j=1}^m S_{p_j}
                          \cup\{c_1,\ldots,c_\ell\}\right),
 \label{eq:deleted-domain}
\end{equation}
the Fueter image in the axial monogenic class has real codimension $8m$,
independent of $\ell$. In particular, every axial monogenic function on
$\HH$ minus finitely many real points has a global slice-regular Fueter
primitive there.
\end{corollary}
\begin{proof}
The symmetric planar domain is obtained by deleting the conjugate pairs
$p_j,\bar p_j$ and the real points $c_k$. Its upper part is just the
upper half-plane with the $m$ points $p_j$ deleted. Small circles around
those points give a winding basis. Apply Theorem~\ref{thm:realaxis}.
\end{proof}

Real punctures are on the boundary of the upper parameter domain, not
holes inside it. This is why they create no affine-monodromy obstruction.
The statement does \emph{not} make their singularities removable and does
not assert that real singularities have no principal parts. For example,
$q^{-1}$ is a global slice-regular primitive, up to normalization, of a
nonremovable real-point Fueter singularity.

\section{A holomorphic invariant and two complex moments}
\label{sec:holomorphic}
The period criterion involves only real first-order data. There is also a
holomorphic encoding that can be especially convenient for rational or
meromorphic examples.

\begin{theorem}[The target-derived second derivative]\label{thm:H}
For $g=P+IQ\in\AM(\Omega_U)$, define
\begin{equation}
 H_g(z)=-\frac12\bigl(P(x,r)+rP_r(x,r)\bigr)
        -\frac{\iota r}{2}P_x(x,r).\label{eq:H}
\end{equation}
This is a single-valued holomorphic $\HC$-valued function on $U$.
Every local holomorphic stem $F$ of a Fueter primitive satisfies
\begin{equation}
 F''=H_g.\label{eq:H-second}
\end{equation}
For each closed curve $\gamma$ in $U$,
\begin{equation}
 a_\gamma(g)=\int_\gamma H_g(z)\dd z,\qquad
 b_\gamma(g)=-\int_\gamma zH_g(z)\dd z.\label{eq:complex-periods}
\end{equation}
In particular, both displayed complex quaternionic integrals actually
belong to the real subspace $\HH\subset\HC$.
\end{theorem}
\begin{proof}
Take a local primitive and its potentials $C,E$ from
Lemma~\ref{lem:forms}. With $F=A+\iota B$ and $F'=\partial_xF$,
\begin{equation}
 F'=C+M,\qquad M=\frac r2Q-\frac{\iota r}{2}P.\label{eq:M}
\end{equation}
Indeed, $A_x=C+rQ/2$ and $B_x=-rP/2$.
Differentiating with respect to $x$ gives
\[
 F''=-P/2+rQ_x/2-\iota rP_x/2=H_g
\]
by $Q_x=-P_r$. This proves local holomorphy, hence global holomorphy
of the explicitly defined function $H_g$. Local primitives differ by
$za+b$, whose second derivative is zero, so \eqref{eq:H-second} has no
choice dependence.

As identities of $\HC$-valued one-forms,
\begin{equation}
 H_g\dd z=\omega_g+dM,\qquad
 zH_g\dd z=d(zM)-\beta_g.\label{eq:exact-corrections}
\end{equation}
The first follows by differentiating \eqref{eq:M}. For the second use
$F=zC+E$ to obtain $zF'-F=zM-E$ and then differentiate. The functions
$M$ and $zM$ are globally single valued. Their differentials integrate to
zero around a closed curve, yielding \eqref{eq:complex-periods}.
The original definitions \eqref{eq:periods} prove the final reality
assertion.
\end{proof}

\begin{corollary}[Two Laurent coefficients at an isolated sphere]
\label{cor:two-residues}
Let $0<R<t$, let $p=s+\iota t$, and let $g$ be axial monogenic on the
circularization of $0<|z-p|<R$. Write the Laurent expansion
\begin{equation}
 H_g(z)=\sum_{n\in\mathbb Z}c_n(z-p)^n,\qquad c_n\in\HC.
 \label{eq:H-laurent}
\end{equation}
Its period pair around a positive small circle is
\begin{equation}
 a=2\pi\iota c_{-1},\qquad
 b=-2\pi\iota\bigl(p c_{-1}+c_{-2}\bigr).\label{eq:residue-pair}
\end{equation}
A global single-valued primitive on that punctured tube exists exactly
when $c_{-1}=c_{-2}=0$. This assertion also holds when $H_g$ has an
essential singularity.
\end{corollary}
\begin{proof}
Apply the ordinary complex residue calculation componentwise to the two
integrals in \eqref{eq:complex-periods}. The punctured disk has one
homology generator, and Theorem~\ref{thm:period} applies. The residue
calculation uses only the Laurent expansion, not meromorphicity.
\end{proof}
A vanishing $c_{-1}$ alone is insufficient. The coefficient $c_{-2}$
retains the constant part of the affine monodromy, even though it is not
the residue of $H_g$ itself. If $H_g$ is meromorphic with only finitely
many isolated poles in a larger simply connected parameter region, the
same criterion is imposed at every deleted point.

For the explicit mode, differentiating \eqref{eq:mode-stem} gives
\begin{equation}
 H_{G_{p;a,b}}=2L_p'a+L_p''(za+b).\label{eq:H-mode}
\end{equation}
Near $p$ its two relevant Laurent coefficients are
\begin{equation}
 c_{-1}=\frac{a}{2\pi\iota},\qquad
 c_{-2}=-\frac{pa+b}{2\pi\iota}.\label{eq:mode-residues}
\end{equation}
Substitution into \eqref{eq:residue-pair} independently checks both the
sign and normalization of \eqref{eq:mode-periods}.

\begin{remark}[The invariant is not injective]\label{rem:H-kernel}
The statement $H_g=0$ does not imply $g=0$ on a product domain.
In fact, \eqref{eq:H} gives $P_x=0$ and $P+rP_r=0$, and the Vekua
system then shows
\begin{equation}
 P=\frac c r,\qquad Q=\frac{cx+d}{r^2},\qquad c,d\in\HH.
 \label{eq:H-kernel}
\end{equation}
These constants are global on connected $U$: $rP$ has both first
partials zero, and $r^2Q-xc$ does as well. Conversely these functions
satisfy \eqref{eq:vekua} and have $H_g=0$.
They are Fueter images of the complex-affine stems
$F=-\tfrac{\iota}{2}(zc+d)$. There is therefore no conflict with the
period criterion: their periods vanish. Smooth continuation through a
real interval would force $c=d=0$.
\end{remark}

\section{Critical spherical singularities: a complete classification}
\label{sec:critical}
Fix $p=s+\iota t$, $t>0$, and choose $0<R<t$. The tubular neighborhood
of the conjugacy sphere $S_p=s+t\SSS$ is
\begin{equation}
 T_R(S_p)=\{x+Ir:(x-s)^2+(r-t)^2<R^2,\ I\in\SSS\}.
 \label{eq:tube}
\end{equation}
It is a genuine Euclidean distance tube: since $r>0$ here,
\begin{equation}
 \dist(x+Ir,S_p)=\sqrt{(x-s)^2+(r-t)^2}=|z-p|.
 \label{eq:distance-tube}
\end{equation}
For a function on the punctured tube define
\begin{equation}
 M_g(\rho)=\sup_{|z-p|=\rho,\ I\in\SSS}|P(z)+IQ(z)|,
 \qquad 0<\rho<R.\label{eq:M-growth}
\end{equation}
Its spherical period pair is
\begin{equation}
 \operatorname{res}_{S_p}(g)=(a,b)
 =\left(\int_{|z-p|=\rho}\omega_g,
        \int_{|z-p|=\rho}\beta_g\right),\label{eq:spherical-residue}
\end{equation}
with positive orientation. Closedness makes this pair independent of
$\rho$. This notation refers to an ordered pair of quaternionic numbers,
not to the flux residue of an isolated point in four dimensions.

\begin{lemma}[A logarithmic bound for an exact target]
\label{lem:log-bound}
Suppose $h\in\AM(T_R(S_p)\setminus S_p)$ has zero spherical period pair
and satisfies $M_h(\rho)=O(\rho^{-1})$. Then every global Fueter
primitive stem $F$ on $0<|z-p|<R$ satisfies
\begin{equation}
 |F(z)|_{\HC}=O\bigl(1+|\log|z-p||\bigr)\quad(z\to p),
 \label{eq:log-bound}
\end{equation}
uniformly in the angular variable.
\end{lemma}
\begin{proof}
By Theorem~\ref{thm:period}, global potentials $C,E$ exist.
Fix a reference circle $|z-p|=\rho_0<R$. For any point of a smaller
circle, take a path along the reference circle to the same angle and
then a radial segment inward. The values and integrals on the fixed
circle are uniformly bounded. Along the radial segment,
\eqref{eq:norm-equivalence} and the growth hypothesis give
$|\omega_h|=O(\rho^{-1})$ and $|\beta_h|=O(\rho^{-1})$, since
$|z|$ is bounded in the disk. Hence both $C$ and $E$ are uniformly
$O(1+|\log\rho|)$. Since $F=zC+E$ and $z$ is bounded, the same is true
of $F$. Adding an affine stem does not change this estimate.
\end{proof}

\begin{lemma}[A holomorphic logarithmic bound is removable]
\label{lem:hol-removable}
If a holomorphic $\HC$-valued function on a punctured disk satisfies
\eqref{eq:log-bound}, it extends holomorphically across the puncture.
\end{lemma}
\begin{proof}
Take its componentwise Laurent series. For every integer $n\ge1$,
the norm of the coefficient of $(z-p)^{-n}$ is at most
$C\rho^n(1+|\log\rho|)$ by the complex coefficient formula.
Letting $\rho\downarrow0$ makes that coefficient zero. All negative
coefficients vanish, so the Laurent series is a Taylor series.
\end{proof}

\begin{theorem}[Complete critical-order spherical principal part]
\label{thm:critical}
Let $g\in\AM(T_R(S_p)\setminus S_p)$ and suppose
\begin{equation}
 M_g(\rho)=O(\rho^{-1})\qquad(\rho\downarrow0).
 \label{eq:critical-bound}
\end{equation}
There are unique $a,b\in\HH$ and a unique $h\in\AM(T_R(S_p))$ such that
\begin{equation}
 g(q)=G_{p;a,b}(q)+h(q)\qquad(q\in T_R(S_p)\setminus S_p).
 \label{eq:critical-decomposition}
\end{equation}
The coefficients are exactly $\operatorname{res}_{S_p}(g)=(a,b)$.
In particular the space of critical-order singular germs modulo regular
germs has quaternionic dimension two, or real dimension eight.
\end{theorem}
\begin{proof}
First, every $G_{p;a,b}$ satisfies \eqref{eq:critical-bound}.
A local logarithm near $p$ is $O(1+|\log\rho|)$ with an angular
argument bounded on a slit disk, while its derivative is $O(\rho^{-1})$.
In terms of its local stem, \eqref{eq:fueter} can be written
\begin{equation}
 P=-\frac2r\ImC F',\qquad
 Q=\frac2r\ReC F'-\frac2{r^2}\ImC F.\label{eq:fueter-stem}
\end{equation}
The quantity $r$ stays bounded away from zero near $p$, so these
estimates prove the required growth uniformly around the circle;
branch independence allows the slit to be placed arbitrarily.
The explicit formula \eqref{eq:mode-P}--\eqref{eq:mode-Q} gives the
same conclusion without branches.

Let $(a,b)$ be the periods of $g$. Subtract the corresponding mode:
$h=g-G_{p;a,b}$. Its periods vanish by
Proposition~\ref{prop:modes}, and it still satisfies the critical growth
bound. By Theorem~\ref{thm:period}, it has a global slice-regular
primitive $f$ on the punctured tube. Lemma~\ref{lem:log-bound} gives a
logarithmic bound for its holomorphic stem. Lemma~\ref{lem:hol-removable}
extends that stem through $p$; reflecting to the lower disk gives a
slice-regular extension through the entire sphere. Its Laplacian is an
axial monogenic extension of $h$ to $T_R(S_p)$.

A regular axial function on the full tube has zero periods, because
the forms extend as closed forms to a disk. Applying periods to any
putative decomposition therefore forces the coefficients to be $(a,b)$.
The remainder is then determined off the sphere and on it by continuity.
Finally every pair occurs by Proposition~\ref{prop:modes}, proving the
dimension statement for germs.
\end{proof}

\begin{corollary}[Sharp subcritical removability]\label{cor:sharp-removal}
An axial monogenic $g$ on the punctured tube satisfying
\begin{equation}
 M_g(\rho)=o(\rho^{-1})\label{eq:subcritical}
\end{equation}
extends uniquely through $S_p$. The little-$o$ cannot be replaced by
big-$O$ for all such functions.
\end{corollary}
\begin{proof}
The length of $|z-p|=\rho$ is $2\pi\rho$, while $|z|$ remains bounded.
The bounds \eqref{eq:period-bounds} show that both periods tend to zero
under \eqref{eq:subcritical}. They are independent of $\rho$, hence
are zero. Theorem~\ref{thm:critical} supplies the extension.
A mode with $(a,b)\ne(0,0)$ has critical big-$O$ growth but nonzero
periods, so it cannot extend. Uniqueness of the extension is continuity.
\end{proof}

The argument proves more than removability below the threshold. At the
threshold itself, it lists \emph{all} possible singular parts, with no
unclassified logarithmic remainder. The logarithms occur in local
primitives; after subtracting the two allowed modes, a single-valued
logarithmically bounded holomorphic primitive must be removable.

\subsection{A constrained Laurent pair, not two arbitrary complex residues}
There is useful extra information hidden in the reality assertion of
Theorem~\ref{thm:H}. Although the two quaternionic periods are independent,
the two complex quaternionic Laurent coefficients are not.

\begin{proposition}[Reality constraint and one-coefficient encoding]
\label{prop:Laurent-reality}
For an arbitrary axial monogenic function on a punctured tube, let
$c_{-1},c_{-2}$ be as in \eqref{eq:H-laurent}. Then
\begin{equation}
 c_{-1}=\frac{\iota}{t}\ReC c_{-2}.\label{eq:Laurent-reality}
\end{equation}
In particular, writing $c_{-2}=u+\iota v$ with $u,v\in\HH$,
\begin{equation}
 a=-\frac{2\pi}{t}u,\qquad
 b=2\pi v+\frac{2\pi s}{t}u.\label{eq:one-coeff}
\end{equation}
A single-valued primitive on the punctured tube therefore exists if and
only if $c_{-2}=0$. The holomorphic invariant $H_g$ cannot have an
isolated pole of order exactly one at $p$.
\end{proposition}
\begin{proof}
By \eqref{eq:residue-pair}, $a,b\in\HH$ and
$c_{-1}=a/(2\pi\iota)$,
$c_{-2}=-(pa+b)/(2\pi\iota)$. Since $p=s+\iota t$,
\[
 c_{-2}=-\frac{t}{2\pi}a+\frac{\iota}{2\pi}(sa+b).
\]
Taking its two central components gives \eqref{eq:one-coeff} and then
\eqref{eq:Laurent-reality}. Thus $c_{-2}=0$ forces $a=b=0$ and also
$c_{-1}=0$. Corollary~\ref{cor:two-residues} proves the criterion.
An isolated simple pole would have $c_{-2}=0$ and $c_{-1}\ne0$,
contradicting \eqref{eq:Laurent-reality}.
\end{proof}
This refinement is specific to an isolated nonreal parameter point.
It does not make the two directly defined one-forms redundant: an arbitrary
loop still has two independent quaternionic periods, and checking the
ordinary residue $c_{-1}$ alone still fails. The single coefficient
$c_{-2}\in\HC$ has exactly the same eight real coordinates as the pair
$(a,b)\in\HH^2$.

\begin{corollary}[Double-pole characterization of critical growth]
\label{cor:double-pole}
For $g\in\AM(T_R(S_p)\setminus S_p)$, the following are equivalent:
\begin{enumerate}[label=\textup{(\roman*)}]
\item $M_g(\rho)=O(\rho^{-1})$.
\item $H_g$ is meromorphic at $p$ with pole order at most two.
\item $g=G_{p;a,b}+h$ with $h$ regular across the sphere.
\end{enumerate}
In this class, nonremovability is equivalent to a pole of order exactly
two of $H_g$. More generally, without any growth hypothesis, $g$ extends
regularly through $S_p$ if and only if $H_g$ extends holomorphically through
$p$.
\end{corollary}
\begin{proof}
Theorem~\ref{thm:critical} gives (i)$\Rightarrow$(iii), and the mode
estimate in its proof gives (iii)$\Rightarrow$(i).
Formula \eqref{eq:H-mode} gives (iii)$\Rightarrow$(ii).
If (ii) holds, subtract the mode corresponding to the two periods.
Equations \eqref{eq:mode-residues} cancel its entire Laurent principal
part, so the remainder $h$ has zero periods and $H_h$ extends
holomorphically. A global primitive stem $F$ of $h$ exists on the
punctured disk. Since $F''=H_h$ extends, its negative Laurent
coefficients all vanish: twice differentiating a term
$d_{-n}(z-p)^{-n}$ gives the nonzero multiplier $n(n+1)$ times
$d_{-n}(z-p)^{-n-2}$. Thus $F$ extends, and its Fueter image extends.
This proves (ii)$\Rightarrow$(iii).

The same last argument applies whenever $H_g$ is holomorphic, since
its two periods then vanish. The reverse implication follows from
\eqref{eq:H}, as $r$ stays positive on the tube. Finally
Proposition~\ref{prop:Laurent-reality} excludes an isolated simple pole,
and \eqref{eq:one-coeff} shows that a nonzero mode has $c_{-2}\ne0$.
\end{proof}

\section{Exact leading norms and logarithmic energy}
\label{sec:energy}
The classification determines whether a singularity is present. We next
determine its exact size. This is the part of the argument in which the
sphere of quaternionic imaginary units must be treated explicitly rather
than replaced by one complex section.

\begin{lemma}[Norms over the imaginary-unit sphere]\label{lem:sphere-norm}
For $P,Q\in\HH$, with ordinary surface measure $d\sigma$ on the unit
sphere $\SSS$,
\begin{align}
 \frac1{4\pi}\int_{\SSS}|P+IQ|^2\dd\sigma(I)
       &=|P|^2+|Q|^2,\label{eq:sphere-mean}\\
 \max_{I\in\SSS}|P+IQ|^2
       &=|P|^2+|Q|^2+2|\ImH(P\bar Q)|.\label{eq:sphere-max}
\end{align}
If $(X,Y)$ is replaced by
$(\alpha X-\beta Y,\beta X+\alpha Y)$, then
\begin{equation}
 \ImH\bigl((\alpha X-\beta Y)
              \overline{(\beta X+\alpha Y)}\bigr)
 =(\alpha^2+\beta^2)\ImH(X\bar Y).
 \label{eq:rotation-cross}
\end{equation}
\end{lemma}
\begin{proof}
Expansion gives
\[
 |P+IQ|^2=|P|^2+|Q|^2-2\ReH(P\bar Q I).
\]
The last term is twice the Euclidean dot product of $\ImH(P\bar Q)$
with $I$. Its spherical mean is zero and its maximum is twice the norm
of that vector, proving the first two formulas.
For \eqref{eq:rotation-cross}, expand the product. The terms
$X\bar X$ and $Y\bar Y$ are real, while
$\ImH(Y\bar X)=-\ImH(X\bar Y)$. The stated identity follows.
\end{proof}

For a residue pair define
\begin{equation}
 c=sa+b,\qquad Z=pa+b=c+\iota t a,\qquad
 \mathcal R_p(a,b)^2=|Z|_{\HC}^2=|c|^2+t^2|a|^2.
 \label{eq:residue-norm}
\end{equation}
This is positive definite in $(a,b)$ because $t>0$.
Under a real translation of the coordinate origin, $b$ changes by the
corresponding real multiple of $a$, but $sa+b$ does not change. Thus
\eqref{eq:residue-norm} is not an artifact of the chosen real origin.

\begin{theorem}[Uniform leading size and exact supremum]
\label{thm:leading}
Let $g$ satisfy the critical growth hypothesis of
Theorem~\ref{thm:critical}, with residue pair $(a,b)$.
Uniformly for $z=p+\rho e^{\iota\theta}$ as $\rho\downarrow0$,
\begin{equation}
 \rho\sqrt{|P(z)|^2+|Q(z)|^2}
 \longrightarrow\frac{\mathcal R_p(a,b)}{\pi t}.\label{eq:axial-leading}
\end{equation}
Furthermore,
\begin{equation}
 \lim_{\rho\downarrow0}\rho M_g(\rho)
 =\frac1{\pi t}
   \sqrt{\mathcal R_p(a,b)^2+2t|\ImH(b\bar a)|}.
 \label{eq:exact-sup}
\end{equation}
In particular, every nonzero residue pair has a strictly positive
critical growth coefficient.
\end{theorem}
\begin{proof}
The regular remainder in \eqref{eq:critical-decomposition} is bounded
on every smaller tube, so it does not affect the limits. Put
$w=z-p=\rho e^{\iota\theta}$. A local branch of the mode stem satisfies
\begin{equation}
 F'_{p;a,b}(z)=\frac{Z}{2\pi\iota w}
                   +O(1+|\log\rho|),\qquad
 F_{p;a,b}(z)=O(1+|\log\rho|),\label{eq:stem-leading}
\end{equation}
in $\HC$ norm, with uniform angular bounds as explained in the proof
of Theorem~\ref{thm:critical}. Write
\[
 K=\frac{Z}{2\pi\iota w}=X+\iota Y.
\]
Since $r=t+O(\rho)$, \eqref{eq:fueter-stem} gives
\begin{equation}
 P=-\frac2tY+O(1+|\log\rho|),\qquad
 Q=\frac2tX+O(1+|\log\rho|).\label{eq:PQ-leading}
\end{equation}
Multiplication by $(2\pi\iota w)^{-1}$ rotates and scales the pair of
central components of $Z$, whence
\[
 |X|^2+|Y|^2=\frac{|Z|_{\HC}^2}{4\pi^2\rho^2}.
\]
This and \eqref{eq:PQ-leading} prove \eqref{eq:axial-leading} uniformly.

To calculate the supremum, apply \eqref{eq:rotation-cross} to the
central components $(c,ta)$ of $Z$. It gives
\[
 |\ImH(X\bar Y)|
  =\frac{t}{4\pi^2\rho^2}|\ImH(c\bar a)|
  =\frac{t}{4\pi^2\rho^2}|\ImH(b\bar a)|,
\]
since $s a\bar a$ is real. For the leading pair
$P_0=-2Y/t$, $Q_0=2X/t$, conjugation reverses the sign of the imaginary
cross component but not its norm. Therefore
\[
 |P_0|^2+|Q_0|^2+2|\ImH(P_0\overline{Q_0})|
 =\frac{\mathcal R_p(a,b)^2+2t|\ImH(b\bar a)|}
        {\pi^2t^2\rho^2}.
\]
This quantity is independent of $\theta$. Formula
\eqref{eq:sphere-max}, together with the uniform
$O(1+|\log\rho|)$ error, proves \eqref{eq:exact-sup}.
The zero-residue case is included, because then $g$ is bounded.
\end{proof}

\begin{remark}[A genuinely quaternionic correction]
The additional term in \eqref{eq:exact-sup} measures the mismatch between
a sphere average and the largest value on one sphere. It is not justified
to replace the four-dimensional supremum by evaluation on one preselected
complex section. Even when $a,b$ commute in a complex plane, the term can
be nonzero; its vanishing condition is specifically
$\ImH(b\bar a)=0$, not merely $ab=ba$.
\end{remark}

\begin{theorem}[Exact tube energy and finite renormalized remainder]
\label{thm:energy}
Under the hypotheses of Theorem~\ref{thm:leading},
\begin{equation}
 \lim_{\rho\downarrow0}
       \rho\int_{\partial T_\rho(S_p)}|g(q)|^2\dd S_q
       =8\mathcal R_p(a,b)^2.\label{eq:energy-trace}
\end{equation}
For every fixed $0<\rho_0<R$, there is a finite real constant
$E_{\mathrm{ren}}=E_{\mathrm{ren}}(g,p,\rho_0)$ such that
\begin{equation}
 \int_{T_{\rho_0}(S_p)\setminus\overline{T_\rho(S_p)}}|g(q)|^2\dd V_q
 =8\mathcal R_p(a,b)^2\log\frac{\rho_0}{\rho}
       +E_{\mathrm{ren}}+o(1).\label{eq:energy-law}
\end{equation}
The critical residue is zero if and only if $g$ is locally square
integrable near the sphere. Equivalently, using the coefficient in
\eqref{eq:H-laurent}, the logarithmic energy coefficient is
\begin{equation}
 8\mathcal R_p(a,b)^2=32\pi^2|c_{-2}|_{\HC}^2.
 \label{eq:energy-Laurent}
\end{equation}
\end{theorem}
\begin{proof}
Parameterize the tube boundary by
\[
 q=s+\rho\cos\theta+I(t+\rho\sin\theta),
 \qquad 0\le\theta<2\pi,
\]
and put $r=t+\rho\sin\theta$. The meridional tangent has length
$\rho$, and each of the two spherical tangents is scaled by $r$.
The corresponding directions are orthogonal. Hence
\begin{equation}
 dS=\rho r^2\dd\theta\,d\sigma(I),\qquad
 dV=\rho r^2\dd\rho\,\dd\theta\,d\sigma(I).
 \label{eq:tube-measures}
\end{equation}
By \eqref{eq:sphere-mean}, the boundary integral is exactly
\begin{equation}
 J(\rho):=\int_{\partial T_\rho}|g|^2\dd S
 =4\pi\rho\int_0^{2\pi}r^2\bigl(|P|^2+|Q|^2\bigr)\dd\theta.
 \label{eq:J}
\end{equation}
The leading estimate \eqref{eq:PQ-leading} gives, uniformly in $\theta$,
\[
 |P|^2+|Q|^2
 =\frac{\mathcal R_p(a,b)^2}{\pi^2t^2\rho^2}
 +O\left(\frac{1+|\log\rho|}{\rho}
                  +(1+|\log\rho|)^2\right).
\]
Since $r^2=t^2+O(\rho)$, substitution into \eqref{eq:J} yields
\begin{equation}
 J(\rho)=\frac{8\mathcal R_p(a,b)^2}{\rho}
       +O\bigl(1+|\log\rho|+\rho(1+|\log\rho|)^2\bigr).
 \label{eq:J-error}
\end{equation}
This proves \eqref{eq:energy-trace}. The error on the right of
\eqref{eq:J-error} is integrable at zero. Integrating in $\rho$, using
\eqref{eq:tube-measures}, proves \eqref{eq:energy-law} with
\[
 E_{\mathrm{ren}}=
 \int_0^{\rho_0}\left(J(u)-\frac{8\mathcal R_p(a,b)^2}{u}\right)\dd u.
\]
For zero residue the function extends smoothly and is locally square
integrable. For nonzero residue the positive coefficient in
\eqref{eq:energy-law} makes the integral diverge.
Finally \eqref{eq:mode-residues}, which also holds for the residue pair
of $g$, gives $Z=-2\pi\iota c_{-2}$. Taking norms proves
\eqref{eq:energy-Laurent}.
\end{proof}

\begin{corollary}[An energy lower bound without a growth assumption]
\label{cor:energy-lower-general}
Let $g$ be any smooth axial monogenic function on a punctured tube, with
period pair $(a,b)$. For $0<\rho<t$ sufficiently small,
\begin{equation}
 \int_{\partial T_\rho}|g|^2\dd S
 \ge \frac{8(t-\rho)^2}{\rho}
       \max\left\{|a|^2,\frac{|b|^2}{(|p|+\rho)^2}\right\}.
 \label{eq:energy-lower-general}
\end{equation}
A nonzero period pair therefore forces logarithmic or faster divergence
of the tube energy even outside the critical growth class.
\end{corollary}
\begin{proof}
On the circle, Cauchy--Schwarz for the first period gives
\[
 |a|^2\le\frac{2\pi\rho}{4}
                 \int_{|z-p|=\rho}(|P|^2+|Q|^2)\dd\ell.
\]
For the second period the same right side is multiplied by at most
$(|p|+\rho)^2$. Formula \eqref{eq:J} also reads
$J(\rho)=4\pi\int_{|z-p|=\rho}r^2(|P|^2+|Q|^2)\dd\ell$,
and $r\ge t-\rho$. Combining these statements proves
\eqref{eq:energy-lower-general}. Integrating its positive multiple of
$1/\rho$ proves the final assertion.
\end{proof}
This last bound is not asserted to have the sharp coefficient of
\eqref{eq:energy-law}. It has a different purpose: it is a purely
period-based obstruction to finite energy that needs no assumed
pointwise growth rate.

\section{Finite spherical principal parts and far-field moments}
\label{sec:principal-parts}
The local classification has a global consequence analogous in spirit to a
finite Mittag--Leffler decomposition, but with critical spherical modes in
place of point poles. We first record the behavior of the modes at infinity.

\begin{proposition}[Exterior expansion of the modes]\label{prop:exterior}
For $|q|>|p|$, the mode has the convergent expansion
\begin{equation}
 G_{p;a,b}(q)=\sum_{n\ge1}\Fu(q^{-n})\,d_n(p;a,b),\label{eq:mode-exterior}
\end{equation}
where
\begin{equation}
 d_n(p;a,b)=-\frac1\pi\left(
     \frac{\ImC(p^{n+1})}{n+1}a
       +\frac{\ImC(p^n)}{n}b\right).\label{eq:mode-dn}
\end{equation}
Here $\ImC(p^n)$ is an ordinary real scalar. The series and its
derivatives converge uniformly on compact subsets of the exterior domain.
In particular $G_{p;a,b}(q)=O(|q|^{-3})$ at infinity.
\end{proposition}
\begin{proof}
On the complex exterior disk, choose the branch near infinity for which
\begin{align*}
 L_p(z)&=\frac1{2\pi\iota}
   \left(\log(1-p/z)-\log(1-\bar p/z)\right)\\
 &=-\sum_{n\ge1}\frac{\ImC(p^n)}{\pi n}\,z^{-n}.
\end{align*}
The coefficients are real. Multiplying by $za+b$ gives a real-affine
constant term $-ta/\pi$ followed by the Laurent coefficients
\eqref{eq:mode-dn}. The convergent exterior stem defines a slice-regular
function on the quaternionic exterior. Applying $\Fu$ kills the
constant and gives \eqref{eq:mode-exterior}. Uniform convergence of
all derivatives on compact exterior subdomains follows from ordinary
holomorphic Laurent convergence, the stem construction, and real
analyticity at the real axis.

Each $\Fu(q^{-n})$ is homogeneous of degree $-n-2$ and smooth on the
unit three-sphere. Scaling the uniformly convergent series from a fixed
radius larger than $|p|$ shows that its leading possible order is
$|q|^{-3}$.
\end{proof}

\begin{lemma}[Entire axial monogenic polynomial growth]
\label{lem:polynomial}
An entire axial monogenic function satisfying
$|h(q)|\le C(1+|q|)^\alpha$, $\alpha\ge0$, is a polynomial in the four
real coordinates of degree at most $N=\lfloor\alpha\rfloor$.
There are unique $c_0,\ldots,c_N\in\HH$ such that
\begin{equation}
 h(q)=\sum_{k=0}^N\mathcal P_k(q)c_k,
 \qquad
 \mathcal P_k(q)=-\frac{\Fu(q^{k+2})}{2(k+1)(k+2)}.
 \label{eq:axial-polynomials}
\end{equation}
The normalization is $\mathcal P_k(x)=x^k$ on $\R$.
\end{lemma}
\begin{proof}
Quaternionic anticommutation gives
$\overline{\Dir}\Dir=\Fu$, so all four real components of $h$ are
harmonic. For a multi-index $\nu$, the ordinary harmonic interior
estimate at a point $q_0$ gives
\[
 |\partial^\nu h(q_0)|\le C_\nu R^{-|\nu|}
                  \sup_{B(q_0,R)}|h|.
\]
For completeness, the constant is obtained by differentiating the
Poisson formula on a unit ball strictly inside a ball of radius two;
rescaling produces $R^{-|\nu|}$. Its value does not depend on $h,R$.
Let $R\to\infty$. Every derivative with $|\nu|>\alpha$ vanishes,
so harmonic real analyticity makes $h$ a polynomial of the stated degree.

There is an entire slice-regular primitive $f$ by
Corollary~\ref{cor:deleted} with no deleted points. Its entire stem $F$
takes values in $\HH$ on the real axis. Taylor expansion at a real $x$ gives
$A(x,r)=F(x)-r^2F''(x)/2+O(r^4)$ and
$B(x,r)=rF'(x)-r^3F'''(x)/6+O(r^5)$.
Equation \eqref{eq:fueter} therefore yields
\begin{equation}
 h(x)=\Fu f(x)=-2F''(x).\label{eq:real-trace}
\end{equation}
Write the real-axis polynomial as $h(x)=\sum_{k=0}^Nx^kc_k$.
The holomorphic identity theorem, applied to the four complex
coordinates, implies $F''(z)=-\tfrac12\sum z^kc_k$ on the whole plane.
Integrating twice and applying $\Fu$ proves
\eqref{eq:axial-polynomials}. The same real-axis calculation gives
$\mathcal P_k(x)=x^k$, and hence uniqueness.
\end{proof}
The polynomial-growth conclusion is standard harmonic analysis; the
particular normalized basis is recorded to make the next classification
fully explicit, not as a new family of polynomials.

\begin{theorem}[Global finite critical principal-part theorem]
\label{thm:finite-parts}
Let $p_1,\ldots,p_m$ be distinct upper-half-plane points and
\[
 \Omega=\HH\setminus\bigcup_{j=1}^mS_{p_j}.
\]
Suppose $g\in\AM(\Omega)$ has critical growth $O(\dist(q,S_{p_j})^{-1})$
near every removed sphere and satisfies
$|g(q)|=O((1+|q|)^\alpha)$ at infinity, for $\alpha\ge0$.
Then, with $N=\lfloor\alpha\rfloor$, there is a unique expansion
\begin{equation}
 g(q)=\sum_{k=0}^N\mathcal P_k(q)c_k
                    +\sum_{j=1}^mG_{p_j;a_j,b_j}(q).
 \label{eq:finite-parts}
\end{equation}
The pairs $(a_j,b_j)$ are exactly the spherical periods of $g$.
Conversely every expression \eqref{eq:finite-parts} has these
regularity and growth properties. The corresponding vector space has
quaternionic dimension $N+1+2m$.

If instead $g(q)\to0$ as $|q|\to\infty$, the polynomial term is absent.
In particular, in this decaying class the $2m$ quaternionic spherical
periods determine the whole function uniquely.
\end{theorem}
\begin{proof}
Subtract the sum of modes with the periods of $g$. Near each sphere,
Theorem~\ref{thm:critical} removes its singularity; all other modes are
smooth there. The remainder extends to an entire axial monogenic
function. Proposition~\ref{prop:exterior} shows that it retains the
specified polynomial growth. Lemma~\ref{lem:polynomial} gives the
polynomial term and its unique coefficients.

The periods recover all mode coefficients, and the real-axis trace of
the resulting entire remainder recovers the polynomial coefficients.
This proves uniqueness and the dimension count. The converse follows
from Lemma~\ref{lem:axial}, the critical mode bound, and the exterior
estimate. If $g\to0$, the entire remainder is bounded and tends to zero;
Lemma~\ref{lem:polynomial} makes it a constant, which must be zero.
\end{proof}
Real punctures are deliberately excluded from this theorem. They have
no additional \emph{inversion obstruction} by Corollary~\ref{cor:deleted},
but their independent singular principal parts would have to be included
in a theorem that also removed real points.

\begin{corollary}[Exact moment conditions for faster far-field decay]
\label{cor:moments}
In the decaying class of Theorem~\ref{thm:finite-parts}, set
\begin{equation}
 D_n=-\frac1\pi\sum_{j=1}^m\left(
  \frac{\ImC(p_j^{n+1})}{n+1}a_j
  +\frac{\ImC(p_j^n)}{n}b_j\right),\qquad n\ge1.
 \label{eq:global-moments}
\end{equation}
For $|q|>\max_j|p_j|$,
\begin{equation}
 g(q)=\sum_{n\ge1}\Fu(q^{-n})D_n.\label{eq:global-exterior}
\end{equation}
For each integer $K\ge1$,
\begin{equation}
 g(q)=O(|q|^{-K-3})\text{ uniformly at infinity}
 \quad\Longleftrightarrow\quad D_1=\cdots=D_K=0.
 \label{eq:moment-equivalence}
\end{equation}
Writing $p_j=s_j+\iota t_j$, the first term is
\begin{equation}
 g(q)=\frac4\pi\frac{\bar q}{|q|^4}
                    \sum_{j=1}^m t_j(s_j a_j+b_j)
                  +O(|q|^{-4}).\label{eq:farfield-first}
\end{equation}
\end{corollary}
\begin{proof}
Sum the finitely many expansions \eqref{eq:mode-exterior}.
If the first $K$ coefficients vanish, homogeneity and convergence give
the indicated uniform decay. Conversely, evaluate the expansion on the
positive real axis. Formula \eqref{eq:real-trace}, applied to $q^{-n}$
on a real neighborhood, gives
\[
 \Fu(q^{-n})\big|_{q=x}=-2n(n+1)x^{-n-2}.
\]
The first nonzero $D_n$ among $n\le K$ would produce a nonzero leading
term with order larger than $x^{-K-3}$, a contradiction. The uniqueness
of Laurent coefficients, or successive multiplication by powers of $x$
and passage to the limit, makes this argument componentwise rigorous.
Finally $D_1=-\pi^{-1}\sum t_j(s_j a_j+b_j)$ and direct differentiation
gives $\Fu(q^{-1})=-4\bar q/|q|^4$, proving
\eqref{eq:farfield-first}.
\end{proof}
A mode may therefore have a nonzero critical spherical residue but a
vanishing first far-field moment. For example $G_{\iota;1,0}$ has
periods $(1,0)$ and decays as $O(|q|^{-4})$, while
$G_{\iota;0,1}$ generally has a nonzero $|q|^{-3}$ term.
A single ``charge at infinity'' cannot distinguish the two residues.

\section{Worked examples and reproducible checks}
\label{sec:examples}
\subsection{One sphere: both obstructions and their energy}
For the unit imaginary sphere, let
\[
 g_0=G_{\iota;0,1},\qquad g_1=G_{\iota;1,0}.
\]
Their local primitives can be taken to be
$\pi^{-1}\arctan z$ and $\pi^{-1}z\arctan z$, respectively.
Their period pairs are $(0,1)$ and $(1,0)$, so neither has a global
single-valued slice-regular primitive on $\HH\setminus\SSS$.
Both have residue norm $\mathcal R=1$. Consequently
\[
 \lim_{\rho\downarrow0}\rho M_{g_j}(\rho)=\frac1\pi,
 \qquad
 \int_{T_{\rho_0}\setminus\overline T_\rho}|g_j|^2\dd V
       =8\log(\rho_0/\rho)+O(1),\qquad j=0,1.
\]
Formula \eqref{eq:mode-real} gives the particularly simple real traces
\[
 g_0(x)=\frac{4x}{\pi(1+x^2)^2},\qquad
 g_1(x)=-\frac4{\pi(1+x^2)^2}.
\]
These traces are nonsingular on $\R$ despite the unavoidable obstruction
around the removed sphere.

For $g_0$, the Laurent residue of $H_{g_0}$ is zero, while its double-pole
coefficient is nonzero:
\[
 c_{-1}=0,\qquad c_{-2}=-\frac1{2\pi\iota}.
\]
This is an explicit counterexample to checking only the ordinary
residue of the target-derived holomorphic invariant.

\subsection{A leading cancellation at one point of a sphere}
Take $p=\iota$, $a=1$, and $b=\ii$, where $\ii$ is a quaternionic unit
and $\iota$ is the central complex unit. The residue norm is
$\mathcal R^2=2$, but $|\ImH(b\bar a)|=1$. Hence
\[
 \lim_{\rho\downarrow0}\rho M_g(\rho)=\frac2\pi,
 \qquad
 \lim_{\rho\downarrow0}\rho\sqrt{|P|^2+|Q|^2}=\frac{\sqrt2}{\pi}.
\]
At $z=\iota+\rho$ the leading term in \eqref{eq:PQ-leading} is
\[
 g(\rho+I)=\frac{\ii+I}{\pi\rho}+O(1+|\log\rho|).
\]
The critical leading term vanishes at $I=-\ii$ and is maximal at
$I=\ii$. Pointwise cancellation along a selected approach to the sphere
therefore does not make the spherical singularity removable.

\subsection{Multiple holes and projection}
On a domain with two holes, choose omitted points $p_1,p_2$ and dual
loops $\gamma_1,\gamma_2$. For instance, a target with period data
\[
 \Per(g)=((\ii,1),(\jj,\kk))
\]
has the unique global decomposition modulo an affine primitive
\[
 g=\Fu f+G_{p_1;\ii,1}+G_{p_2;\jj,\kk}.
\]
Subtracting only one mode per hole cannot remove arbitrary data: each
hole has both a slope and a constant monodromy coordinate. The
projection in \eqref{eq:projection} removes precisely all four
quaternionic coefficients, not an arbitrary fit to the boundary values.

\subsection{Checks included with this article}
The archive contains \texttt{verify\_results.py} and its recorded output.
The program uses NumPy and SymPy and performs the following independent
checks: symbolic substitution in both Vekua equations; symbolic
comparison of \eqref{eq:H} with the second derivative of a local
logarithmic stem; numerical contour integration of both periods; direct
comparison of the branch-free mode formula with
\eqref{eq:fueter-stem}; and checks of the exact norm and energy constants.
The numerical tests include genuinely noncommuting right coefficients.

For a reproducible example, use
\[
 p=\frac3{10}+\frac65\iota,\quad
 a=1-2\ii+\tfrac12\jj+\kk,\quad
 b=\tfrac12+\ii-\jj+2\kk.
\]
Then $\mathcal R_p(a,b)^2=253/16$, so the predicted tube trace limit is
$253/2=126.5$. The predicted supremum coefficient is approximately
$1.47242217559$. Evaluation of the explicit formulas gives:
\begin{center}
\renewcommand{\arraystretch}{1.18}
\begin{tabular}{@{}rcc@{}}
\toprule
$\rho$ & $\rho M_g(\rho)$ &
$\displaystyle\rho\int_{\partial T_\rho}|g|^2\dd S$\\
\midrule
$10^{-2}$ & $1.49117985106$ & $126.669179900$\\
$10^{-3}$ & $1.47469388170$ & $126.503970079$\\
$10^{-4}$ & $1.47269053008$ & $126.500071800$\\
$10^{-5}$ & $1.47245325267$ & $126.500001132$\\
\midrule
Proved limit & $1.47242217559\ldots$ & $126.5$\\
\bottomrule
\end{tabular}
\end{center}
The supremum is evaluated using \eqref{eq:sphere-max}, avoiding a
sampling error over imaginary units; only the meridional angle is
sampled. The spherical mean is exact. These tests detect normalization,
sign, and multiplication-order mistakes. They neither prove the
asymptotic theorem nor establish its novelty; those are separate matters.

\section{Relation to the literature and the scope of the contribution}
\label{sec:position}
\subsection{What is established background}
Inverse Fueter theory was developed by Colombo, Sabadini, and Sommen
\cite{CSS}; their later integral treatment uses spherical monogenics
\cite{CSSIntegral}. The direct integration construction in
\cite[Theorem~3]{CPSS} is stated on a rectangular parameter region.
Perotti's surjectivity theorem \cite[Theorem~2]{Perotti} explicitly assumes
that every connected component of the symmetric planar set is simply
connected. These references support local and appropriately qualified
global inversion; we do not treat that foundational fact as new.

The local affine kernel, the stem description of slice functions, the
axial Vekua equations, and the classical complex operations used here are
also background. In particular, Remark~\ref{rem:known-modes} identifies
our elementary modes with previously studied spherical Cauchy kernels in
the unit-sphere case. Repackaging those kernels without acknowledging
their source would not constitute a new theorem.

The 2026 preprint by De Martino and Pinton \cite{DMP2026} concerns a
variation of inverse Fueter theory and a polyanalytic
Cauchy--Kovalevskaya extension. Its abstract and introductory framework
were included in the current-literature check. The present paper concerns
the ordinary four-dimensional Laplacian, global affine periods, and
critical spherical energy, not that polyanalytic extension problem.

\subsection{The proposed contributions}
The results whose precise formulations were not located in the inspected
primary sources are the following connected package:
\begin{enumerate}[label=\textup{(\alph*)}]
\item the two simultaneously target-defined forms, the explicit complete
monodromy homomorphism, and the finite-topology split cokernel
(Theorems~\ref{thm:period}, \ref{thm:monodromy}, and \ref{thm:split});
\item the target-derived holomorphic invariant, its exact differential
corrections, and the admissibility constraint on its two relevant
Laurent coefficients (Theorem~\ref{thm:H} and
Proposition~\ref{prop:Laurent-reality});
\item the complete critical spherical principal part, the double-pole
characterization, and especially the exact leading supremum and
renormalized energy formulas
(Theorem~\ref{thm:critical}, Corollary~\ref{cor:double-pole}, and
Theorems~\ref{thm:leading}--\ref{thm:energy});
\item the quantitative non-approximation bound, the real-versus-spherical
puncture distinction, and the global finite principal-part and far-field
moment classifications (\S\ref{sec:stability},
Corollary~\ref{cor:deleted}, and \S\ref{sec:principal-parts}).
\end{enumerate}
The global period calculation is elementary once the correct two
potentials have been found. The stronger analytic content is that the
same eight real coordinates exhaust all critical spherical singularities
and determine a precise positive quadratic energy coefficient. That
classification and its quantitative consequences are the main reason the
problem is more substantial than merely observing that closed forms can
have periods.

This is a qualified novelty statement. The literature search is not a
proof that no equivalent theorem exists, particularly in older generalized
analytic-function or Clifford-analysis treatments. The manuscript does
not claim that the inverse Fueter problem as a whole remained unsolved,
nor that it corrects a theorem stated with its proper domain and branch
hypotheses. The attached expository manuscript is used as motivation and
as the source of the two-integration starting point, not as a published
priority reference.

\subsection{Precise limits of what has been proved}
The target class is axial and left monogenic, and primitives are
single-valued slice functions induced by holomorphic stems. The results
do not assert that every nonaxial Fueter-regular function is a Laplacian
of a slice-regular function. Finite-dimensional cokernel statements use
the explicitly stated finite winding-basis hypothesis; the all-loop
criterion itself needs no finite topology. The singularity and exact
energy classification applies to a smooth conjugacy sphere of positive
radius, not to a collapsed real point. The normalized inverse estimates
are compact-local and do not claim uniform conditioning when holes
collide or approach the real axis.

Natural subsequent problems include uniform estimates for degenerating
families of parameter domains and a higher-dimensional version with the
larger polynomial kernel of the Fueter--Sce map. Neither extension is
needed for any theorem proved here, and neither is claimed as completed.

\section*{Conclusion}
\addcontentsline{toc}{section}{Conclusion}
Global Fueter inversion on a punctured or multiply connected parameter
domain is governed by two explicit closed forms. Their periods are both
necessary and sufficient, admit an explicit splitting, and give stable
lower bounds against global approximation. At a nonreal puncture, the
same obstruction becomes a constrained Laurent pair, equivalently one
complex quaternionic double-pole coefficient. At critical spherical
growth, it is not merely an obstruction to a primitive: it is the complete
singular part. Its norm determines exactly the coefficient of logarithmic
four-dimensional energy divergence. This supplies a unified global,
local, and quantitative answer in the axial quaternionic setting.

\begin{thebibliography}{9}
\addcontentsline{toc}{section}{References}

\bibitem{GP}
R.~Ghiloni and A.~Perotti,
\emph{Slice regular functions on real alternative algebras},
Advances in Mathematics \textbf{226} (2011), 1662--1691.
\href{https://arxiv.org/abs/1008.4318}{arXiv:1008.4318}.

\bibitem{CSS}
F.~Colombo, I.~Sabadini, and F.~Sommen,
\emph{The inverse Fueter mapping theorem},
Communications on Pure and Applied Analysis \textbf{10} (2011), 1165--1181.
\href{https://doi.org/10.3934/cpaa.2011.10.1165}{doi:10.3934/cpaa.2011.10.1165}.

\bibitem{CSSIntegral}
F.~Colombo, I.~Sabadini, and F.~Sommen,
\emph{The inverse Fueter mapping theorem in integral form using spherical
monogenics}, Israel Journal of Mathematics \textbf{194} (2013), 485--505.
\href{https://doi.org/10.1007/s11856-012-0090-4}{doi:10.1007/s11856-012-0090-4}.

\bibitem{CPSS}
F.~Colombo, D.~Pe\~na Pe\~na, I.~Sabadini, and F.~Sommen,
\emph{A new integral formula for the inverse Fueter mapping theorem},
Journal of Mathematical Analysis and Applications \textbf{417} (2014),
112--122. \href{https://arxiv.org/abs/1302.0685}{arXiv:1302.0685}.
The theorem and example numbering cited here is that of the arXiv version.

\bibitem{Perotti}
A.~Perotti,
\emph{A Local Cauchy Integral Formula for Slice-Regular Functions},
Computational Methods and Function Theory \textbf{24} (2024), 185--203;
first published online 30 May 2023.
\href{https://doi.org/10.1007/s40315-023-00485-5}{doi:10.1007/s40315-023-00485-5}.

\bibitem{DMP2026}
A.~De Martino and S.~Pinton,
\emph{A variation of the inverse Fueter theorem and the generalized
polyanalytic Cauchy--Kovalevskaya extension of order 2},
preprint (2026).
\href{https://arxiv.org/abs/2608.07711}{arXiv:2608.07711}.

\bibitem{companion}
{\raggedright
\emph{Classical Complex Theorems over the Quaternions:
Slice-regular and Cauchy--Fueter analogues, with proofs},
user-supplied expository manuscript, 20 September 2026,
\texttt{quaternionic\_analysis.tex}, especially the final section on
Fueter's mapping theorem and its constructive local inverse.
Unpublished companion material; not a priority source.\par}

\end{thebibliography}
\end{document}
